\documentclass{article}
\usepackage[T1]{fontenc}
\usepackage{lmodern}
\usepackage{graphicx}
\usepackage{amsmath,amssymb}
\usepackage[a4paper,margin=2cm]{geometry}
\usepackage[absolute,overlay]{textpos}
\setlength{\TPHorizModule}{1mm}
\setlength{\TPVertModule}{1mm}
\usepackage[indonesian]{babel}
\usepackage{hyperref}
\setlength{\emergencystretch}{2em}
% unit: Halaman 35

\begin{document}

% === Halaman 35 (python/native) ===

ExpansionKey()

Algoritma: 

1.

Salin elemen-elemen key ke dalam larik w[0], w[1], w[2], w[3]. Larik w[0] berisi 

empat elemen pertama key, w[1] berisi empat elemen berikutnya, dan seterusnya.

2.

Mulai dari i = 4 sampai 43, lakukan:

a) Simpan w[i-1] ke dalam peubah temp

b) Jika i kelipatan 4, lakukan fungsi g berikut:

• RotWord: Geser w[i-1] satu byte ke kiri secara sirkuler

• SubWord: Lakukan substitusi dengan S-box terhadap hasil pergeseran tersebut

\end{document}
